\chapter{Language implementation}

\section{Lexer Implementation}

The lexer is responsible for transforming raw source text into a flat sequence of tokens, which the parser then consumes. It operates line by line, handling both the structural indentation logic and the inline tokenization of each line's content.

\subsection{Token Recognition}

Token recognition is driven by a list of ordered token specifications, each pairing a \texttt{TokenType} with a regular expression pattern. The order is significant --- more specific or longer patterns are listed before shorter ones to prevent incorrect partial matches. For example, the two-character operators \texttt{!=}, \texttt{==}, \texttt{<=}, and \texttt{>=} are specified before their single-character counterparts \texttt{<} and \texttt{>}.

All individual patterns are combined at class initialization into a single master regular expression using alternation, and compiled once into a \texttt{Pattern} object. During tokenization, a \texttt{Matcher} is advanced through the content of each line using \texttt{lookingAt()}, which attempts to match at the current position without requiring a full-string match. The capturing group that produced the match is then used to identify the corresponding \texttt{TokenSpec}.

\subsection{Indentation Handling}

Since the DSL uses indentation-based block delimitation, the lexer is responsible for emitting \texttt{INDENT} and \texttt{DEDENT} tokens explicitly. 
An integer stack tracks the indentation levels encountered so far, initialized with a base level of zero. At the start of each line, the number of leading spaces and tabs is counted. 
If the measured indentation exceeds the top of the stack, an \texttt{INDENT} token is emitted and the new level is pushed. If it is less, \texttt{DEDENT} tokens are emitted and levels are popped until the stack matches the current indentation. 
Blank and whitespace-only lines are skipped entirely and carry no structural meaning.

\subsection{Keyword Resolution and Special Cases}

Identifiers are first matched by a general alphanumeric pattern, then passed through a keyword lookup to determine whether the matched text corresponds to a reserved word such as \texttt{for}, \texttt{while}, \texttt{if}, or \texttt{true}. If a keyword match is found, the corresponding keyword token type is emitted instead of a generic \texttt{IDENT} token.

String literals have their surrounding quotation marks stripped during tokenization, so the parser receives only the inner content. Comments, introduced by a double dash (\texttt{--}), cause the lexer to discard the remainder of the line immediately upon detection. Any character that does not match any known pattern results in an \texttt{ILLEGAL} token carrying the unrecognized character, allowing the parser to produce a meaningful error message rather than failing silently.

Each line concludes with an explicit \texttt{NEWLINE} token. At the end of the source, any remaining open indentation levels are closed with corresponding \texttt{DEDENT} tokens, followed by a final \texttt{EOF} token.

\section{Parser Implementation}

The parser consumes the flat token sequence produced by the lexer and constructs an Abstract Syntax Tree (AST), which represents the hierarchical structure of the program. It is implemented as a hand-written recursive descent parser, where each grammar rule corresponds directly to a parsing method.

\subsection{Token Navigation}

The parser maintains a cursor over the token list and exposes three core navigation primitives. The \texttt{advance()} method moves the cursor forward and returns the new current token. The \texttt{check()} method tests whether the current token matches a given type without consuming it. The \texttt{expect()} method asserts that the current token has a specific type, consumes it, and throws a \texttt{ParseException} with a descriptive message including the line and column number if the assertion fails. A \texttt{peek()} method allows looking ahead by an arbitrary offset, which is used for disambiguation without side effects.

\subsection{Statement Parsing and Disambiguation}

The top-level entry point, \texttt{parseProgram()}, skips any leading newlines and delegates to \\ \texttt{parseStatementList()}, which iterates until an \texttt{EOF} or, when parsing inside a block, a \texttt{DEDENT} token is encountered.

Individual statements are dispatched in \texttt{parseStatement()} based on the type of the current token. Control flow keywords (\texttt{for}, \texttt{while}, \texttt{if}) are routed directly to their respective parsing methods. When the current token is an identifier, a one-token lookahead is used to disambiguate between an assignment, where the next token is \texttt{=}, and a function call statement, where the next token is \texttt{(}. Any identifier-headed expression that resolves to something other than a function call is rejected with a parse error, enforcing the grammar's restriction that bare expressions are not valid statements.

\subsection{Control Flow Constructs}

The two forms of \texttt{for} loop are disambiguated after consuming the \texttt{for} keyword by checking whether the next token is an identifier followed by the \texttt{from} keyword. If so, a range loop is parsed; otherwise, a count loop is assumed. Both forms conclude by parsing an indented block via \texttt{parseBlock()}, which expects an \texttt{INDENT} token, recursively parses a statement list, and then expects a \texttt{DEDENT} token. The \texttt{if} statement optionally parses an \texttt{else} branch by checking for the \texttt{else} keyword after the then-block, with intervening newlines skipped.

\subsection{Expression and Condition Parsing}

Expressions are parsed using a standard two-level recursive descent that naturally encodes operator precedence. The \texttt{parseExpression()} method handles additive operators by parsing a \texttt{parseTerm()} and then consuming any sequence of \texttt{+} or \texttt{-} operators followed by additional terms, constructing left-associative \texttt{BinaryOp} nodes. The \texttt{parseTerm()} method handles multiplicative operators in the same manner, and \texttt{parseFactor()} handles the base cases: numeric, string, and boolean literals; parenthesised expressions; unary negation; and identifier-headed constructs.

When an identifier is encountered in \texttt{parseFactor()}, a further one-token lookahead determines whether it begins a function call (\texttt{IDENT} followed by \texttt{(}), a dot-notation special object access (\texttt{IDENT} followed by \texttt{.}), or a plain variable reference.

Conditions are parsed separately from expressions. The \texttt{parseCondition()} method first checks for a \texttt{not} keyword, in which case it recursively parses a negated condition. Otherwise, it parses a left-hand expression and checks whether the current token is a comparison operator. If so, a \texttt{ComparisonCondition} node is produced; if not, the expression itself is wrapped in a \texttt{BooleanCondition}, supporting the use of bare identifiers and function calls as boolean predicates.

\subsection{Error Handling}

All syntax errors are reported by throwing a \texttt{ParseException}, an unchecked exception that carries a human-readable message including the offending token type, its literal value where applicable, and its line and column position. This ensures that error messages are actionable for the learner rather than presenting raw internal state.

\section{AST Examples}

The following examples illustrate the simplified abstract syntax trees produced by the parser for various language constructs. Each tree omits certain lexical details (such as spaces and optional formatting) to highlight the essential structure.

\subsection{Variable Assignment}

The assignment statement \texttt{x = 45 * 3 + 2} demonstrates correct operator precedence and nesting of expressions.

\begin{center}
\begin{forest}
for tree={
    font=\small,
    grow=south,
    parent anchor=south,
    child anchor=north,
    l sep=6mm,
    s sep=4mm,
    inner sep=2pt
}
[{$<$assignment$>$}
    [{$<$identifier$>$}[{"x"},tier=word]]
    [{"="}[{"="},tier=word]]
    [{$<$expression$>$}
        [{$<$expression$>$}
            [{$<$term$>$}
                [{$<$term$>$}
                    [{$<$number$>$}
                        [{"45"},tier=word]
                    ]
                ]
                [{"*"},tier=word]
                [{$<$factor$>$}
                    [{"3"},tier=word]
                ]
            ]
        ]
        [{"+"}[{"+"},tier=word]]
        [{$<$term$>$}
            [{"2"},tier=word]
        ]
    ]
]
\end{forest}
\end{center}

\noindent Corresponding source line:
\begin{lstlisting}
x = 45 * 3 + 2
\end{lstlisting}

\subsection{For Count Loop}

The count loop \texttt{for 2 times x = 1} shows how the block structure is captured via \texttt{INDENT} and \texttt{DEDENT} markers, similar to Python.

\begin{center}
\begin{forest}
for tree={
    font=\small,
    grow=south,
    parent anchor=south,
    child anchor=north,
    l sep=6mm,
    s sep=4mm,
    inner sep=2pt
}
[{$<$for\_count$>$}
    [{"for"}[{"for"},tier=word]]
    [{$<$expression$>$}
        [{$<$number$>$}[{"2"},tier=word]]
    ]
    [{"times"}[{"times"},tier=word]]
    [{NEWLINE}
        [{"\textbackslash n"},tier=word]
    ]
    [{$<$block$>$}
        [{INDENT}[{$\rightarrow$},tier=word]]
        [{$<$assignment$>$}
            [{$<$identifier$>$}[{"x"},tier=word]]
            [{"="}[{"="},tier=word]]
            [{$<$expression$>$}
                [{$<$number$>$}[{"1"},tier=word]]
            ]
        ]
        [{DEDENT}[{$\leftarrow$},tier=word]]
    ]
]
\end{forest}
\end{center}

\noindent Corresponding source lines:
\begin{lstlisting}
for 2 times
    x = 1
\end{lstlisting}

\subsection{For Range Loop}

The range loop \texttt{for x from 1 to 9 y = 1} iterates a variable over a sequence of values. The AST explicitly represents the loop variable, the bounds, and the body.

\begin{center}
\begin{forest}
for tree={
    font=\tiny,
    grow=south,
    parent anchor=south,
    child anchor=north,
    l sep=4mm,
    s sep=1mm,
    inner sep=1pt
}
[{$<$for\_range$>$}
    [{"for"}[{"for"},tier=word]]
    [{$<$identifier$>$}[{"x"},tier=word]]
    [{"from"}[{"from"},tier=word]]
    [{$<$expression$>$}
        [{$<$number$>$}[{"1"},tier=word]]
    ]
    [{"to"}[{"to"},tier=word]]
    [{$<$expression$>$}
        [{$<$number$>$}[{"9"},tier=word]]
    ]
    [{NEWLINE}
        [{"\textbackslash n"},tier=word]
    ]
    [{$<$block$>$}
        [{INDENT}[{$\rightarrow$},tier=word]]
        [{$<$assignment$>$}
            [{$<$identifier$>$}[{"y"},tier=word]]
            [{"="}[{"="},tier=word]]
            [{$<$expression$>$}
                [{$<$number$>$}[{"1"},tier=word]]
            ]
        ]
        [{DEDENT}[{$\leftarrow$},tier=word]]
    ]
]
\end{forest}
\end{center}

\noindent Corresponding source lines:
\begin{lstlisting}
for x from 1 to 9
    y = 1
\end{lstlisting}

\subsection{While Loop}

The while loop \texttt{while x > 0 y = 1} demonstrates a condition node containing a comparison operator.

\begin{center}
\begin{forest}
for tree={
    font=\tiny,
    grow=south,
    parent anchor=south,
    child anchor=north,
    l sep=4mm,
    s sep=1mm,
    inner sep=1pt
}
[{$<$while\_loop$>$}
    [{"while"}[{"while"},tier=word]]
    [{$<$condition$>$}
        [{$<$expression$>$}
            [{$<$identifier$>$}[{"x"},tier=word]]
        ]
        [{$<$comparison\_op$>$}[{"$>$"},tier=word]]
        [{$<$expression$>$}
            [{$<$number$>$}[{"0"},tier=word]]
        ]
    ]
    [{NEWLINE}[{"\textbackslash n"},tier=word]]
    [{$<$block$>$}
        [{INDENT}[{$\rightarrow$},tier=word]]
        [{$<$assignment$>$}
            [{$<$identifier$>$}[{"y"},tier=word]]
            [{"="}[{"="},tier=word]]
            [{$<$expression$>$}
                [{$<$number$>$}[{"1"},tier=word]]
            ]
        ]
        [{DEDENT}[{$\leftarrow$},tier=word]]
    ]
]
\end{forest}
\end{center}

\noindent Corresponding source lines:
\begin{lstlisting}
while x > 0
    y = 1
\end{lstlisting}

\subsection{If-Else Statement}

The conditional statement \texttt{if x == 0 y = 1 else z = 2} shows the structure of an if-else with two separate blocks.

\begin{center}
\noindent\makebox[\textwidth][c]{
    \scalebox{0.9}{
\begin{forest}
for tree={
    font=\tiny,
    grow=south,
    parent anchor=south,
    child anchor=north,
    l sep=4mm,
    s sep=1mm,
    inner sep=1pt
}
[{$<$if\_stmt$>$}
    [{"if"}[{"if"},tier=word]]
    [{$<$condition$>$}
        [{$<$expression$>$}
            [{$<$identifier$>$}[{"x"},tier=word]]
        ]
        [{$<$comparison\_op$>$}[{"=="},tier=word]]
        [{$<$expression$>$}
            [{$<$number$>$}[{"0"},tier=word]]
        ]
    ]
    [{NEWLINE}[{"\textbackslash n"},tier=word]]
    [{$<$block$>$}
        [{INDENT}[{$\rightarrow$},tier=word]]
        [{$<$assignment$>$}
            [{$<$identifier$>$}[{"y"},tier=word]]
            [{"="}[{"="},tier=word]]
            [{$<$expression$>$}
                [{$<$number$>$}[{"1"},tier=word]]
            ]
        ]
        [{DEDENT}[{$\leftarrow$},tier=word]]
    ]
    [{"else"}[{"else"},tier=word]]
    [{NEWLINE}[{"\textbackslash n"},tier=word]]
    [{$<$block$>$}
        [{INDENT}[{$\rightarrow$},tier=word]]
        [{$<$assignment$>$}
            [{$<$identifier$>$}[{"z"},tier=word]]
            [{"="}[{"="},tier=word]]
            [{$<$expression$>$}
                [{$<$number$>$}[{"2"},tier=word]]
            ]
        ]
        [{DEDENT}[{$\leftarrow$},tier=word]]
    ]
]
\end{forest}
    }
}
\end{center}

\noindent Corresponding source lines:
\begin{lstlisting}
if x == 0
    y = 1
else
    z = 2
\end{lstlisting}

These examples confirm that the parser correctly handles operator precedence, block structure, and the various control flow constructs defined in the language. The simplified ASTs omit whitespace and delimiters, focusing on the semantic relationships that will be used in later compilation phases.
The full variants of these ASTs are included in the Appendix~\ref{appendix:asts}.

\section{Interpreter Implementation}

The interpreter walks the AST produced by the parser and executes the program, producing a sequence of \texttt{Action} objects that are later dispatched to the Minecraft game engine.
The interpreter is intentionally data-driven, separating core AST evaluation logic from extensible function bindings and special object resolution.

\subsection{Overall Architecture}

The \texttt{Interpreter} class serves as the main entry point for execution.
It maintains references to a \texttt{FunctionRegistry} (for user-defined and standard library functions) and a \texttt{SpecialObjectResolver} (for dot-notation object access), making the system extensible without coupling the traversal logic to specific function implementations.
Upon initialization, an \texttt{Environment} is created to hold the program's variables, and lists are allocated for emitted \texttt{Action} objects and logged messages.

The primary method, \texttt{execute(ASTNode.Program program)}, accepts a parsed AST and returns an \texttt{InterpreterResult} containing the final value, all emitted actions, a snapshot of the final environment state, all logged messages, and a stopped flag indicating whether a \texttt{stop} statement was encountered.

\subsection{Environment and Value Management}

The \texttt{Environment} class manages variable storage using a stack of scopes.
Each scope is a \\ \texttt{LinkedHashMap<String, Value>} mapping variable names to their runtime values.
Although the current language uses only global scope, the stack structure permits future extensions for nested scopes (e.g., function-local variables) without API changes.

Variable assignment is handled by the \texttt{assign(name, value)} method, which searches the scope stack from the topmost scope downward; if the variable already exists in any scope, its value is updated in place.
Otherwise, the variable is created in the topmost (most local) scope.
Variable lookup via \texttt{get(name)} searches the stack from top to bottom and throws an \texttt{InterpreterException} if the variable is undefined.

The \texttt{Value} class is a wrapper around arbitrary \texttt{Object} instances, providing type-aware conversion methods.
A single \texttt{NULL} instance is cached and reused for null values.
The \texttt{asLong()}, \texttt{asBoolean()}, and \texttt{asString()} methods perform implicit type conversions---for example, boolean \texttt{true} converts to \texttt{1L}, and any non-empty string converts to \texttt{true}---allowing programs to mix types flexibly.

\subsection{Statement Execution}

Statements are executed by the \texttt{executeStatement()} method, which dispatches on the AST node type and handles all control flow constructs.

\subsubsection{Assignments and Variables}

An \texttt{Assignment} node is handled by evaluating the right-hand expression and storing the result in the environment under the target variable name.
Assignment Execution (Pseudocode):

\begin{lstlisting}
if (node instanceof ASTNode.Assignment assignment) {
    Value value = evaluateExpression(assignment.value, ...);
    environment.assign(assignment.name, value);
    return value;
}
\end{lstlisting}

\subsubsection{For Loops}

The interpreter handles two forms of \texttt{for} loops.
For count loops (\texttt{for N times}), the count expression is evaluated to a long, and the loop body is executed that many times.
For range loops (\texttt{for X from A to B}), the bounds are evaluated, and the loop variable is assigned to each value from \texttt{A} (inclusive) to \texttt{B} (exclusive) in sequence, allowing the body to access the current iteration value.

Both loop forms catch \texttt{StopSignal} exceptions (thrown by a \texttt{stop} statement) to enable early exit.
When caught, the loop terminates immediately and execution resumes after the loop.
For Count Loop Execution (Pseudocode):

\begin{lstlisting}
if (node instanceof ASTNode.ForCount forCount) {
    long count = evaluateExpression(forCount.count, ...).asLong();
    Value last = Value.nullValue();
    for (long i = 0; i < count; i++) {
        try {
            last = executeStatements(forCount.body.statements, ...);
        } catch (StopSignal stopSignal) {
            break;
        }
    }
    return last;
}
\end{lstlisting}

\subsubsection{While Loops}

While loops evaluate their condition repeatedly using \texttt{evaluateCondition()}, which can handle various condition types (comparison conditions, boolean conditions, and negations).
The loop body is executed as long as the condition evaluates to a truthy value. As with \texttt{for} loops, a \texttt{stop} statement breaks out of the loop.

\subsubsection{If-Else Statements}

The condition is evaluated once, and if it is truthy, the then-block is executed.
If the condition is falsy and an else-block is present, the else-block is executed instead.
Only one branch is ever executed.

\subsubsection{Stop Statements}

When a \texttt{stop} statement is encountered, the interpreter throws an internal \texttt{StopSignal} exception.
This signal propagates up through nested loops and blocks, causing immediate termination.
The top-level \texttt{execute()} method catches this signal and sets the stopped flag in the returned result.

\subsection{Expression Evaluation}

Expressions are evaluated by the \texttt{evaluateExpression()} method, which similarly dispatches on node type.

\subsubsection{Literals}

Numeric, string, and boolean literals are wrapped in \texttt{Value} objects. Identifiers are looked up in the environment.

\subsubsection{Arithmetic and String Operations}

Binary operators are handled by evaluating both operands and applying the operator:
\begin{itemize}
\item[--] Addition (\texttt{+}) performs numeric addition if both operands are numbers, otherwise concatenates them as strings.
\item[--] Subtraction, multiplication, division, and modulo perform numeric operations after converting operands to long.
\item[--] Division and modulo by zero throw an \texttt{InterpreterException}.
\end{itemize}

\subsubsection{Function Calls}

When a \texttt{FunctionCall} node is encountered, all arguments are first evaluated to \texttt{Value} objects.
An \texttt{ExecutionContext} is constructed containing references to the interpreter, environment, Minecraft context (if available), the list of emitted actions, the list of logs, and the evaluated arguments.
The function is then looked up in the \texttt{FunctionRegistry} and invoked with the context.
Function Call Execution (Pseudocode):

\begin{lstlisting}
if (expression instanceof ASTNode.FunctionCall functionCall) {
    List<Value> arguments = new ArrayList<>();
    for (ASTNode argument : functionCall.arguments) {
        arguments.add(evaluateExpression(argument, ...));
    }
    ExecutionContext context = baseContext(..., arguments);
    return functionRegistry.invoke(functionCall.name, context);
}
\end{lstlisting}

\subsubsection{Special Object Access}

Dot notation (e.g., \texttt{player.direction}) is handled by passing the object name and field name to the \texttt{SpecialObjectResolver}, which returns the corresponding value.
This design allows the standard library to register special objects without modifying the interpreter core.

\subsection{Condition Evaluation}

Conditions are parsed into dedicated AST nodes and evaluated by the \texttt{evaluateCondition()} method, which distinguishes several cases:

\begin{itemize}
\item[\textbf{--}] \textbf{Comparison Conditions:} Both operands are evaluated as expressions, then compared using the specified operator. The result is wrapped in a \texttt{Value}.
\item[\textbf{--}] \textbf{Boolean Conditions:} A bare expression (e.g., an identifier or function call) is evaluated, and its truthiness is returned.
\item[\textbf{--}] \textbf{Negation:} A \texttt{not} keyword causes the operand condition to be recursively evaluated and negated.
\item[\textbf{--}] \textbf{Logical Operators:} \texttt{and} and \texttt{or} operators employ short-circuit evaluation---\texttt{and} returns false as soon as the left operand is false, and \texttt{or} returns true as soon as the left operand is true.
\end{itemize}

\subsection{Function Registry and Function Binding}

The \texttt{FunctionRegistry} maintains a map of function names (normalized to strings) to \\ \texttt{InterpreterFunction} instances.
Functions are registered by name during interpreter construction, typically via the \texttt{StandardLibrary} class, which installs all built-in functions.

An \texttt{InterpreterFunction} is a functional interface with a single method \\ \texttt{invoke(ExecutionContext context)} that returns a \texttt{Value}. This design permits functions to:
\begin{itemize}
\item[--] Inspect and modify the environment.
\item[--] Query the current Minecraft state via \texttt{minecraftContext}.
\item[--] Emit \texttt{Action} objects to be executed in sequence.
\item[--] Log messages.
\item[--] Access their arguments via the \texttt{ExecutionContext}.
\end{itemize}

The standard library includes functions for basic I/O (e.g., \texttt{log}, \texttt{print}), player movement (e.g., \texttt{move\_forward}), inventory management (e.g., \texttt{open\_inventory}), and timing (e.g., \texttt{wait}).

\subsection{Error Handling}

All runtime errors are reported via \texttt{InterpreterException}, an unchecked exception that carries a descriptive message.
Type conversion errors, undefined variables, division by zero, and unknown functions all throw this exception with appropriate context.
This approach ensures that parse-time errors are distinct from runtime errors, aiding debugging.

\section{Minecraft Integration and Action Execution}

The CodeCraft DSL is executed within a Minecraft client environment where game logic proceeds at a fixed rate of 20 game ticks per second (50 milliseconds per tick).
The DSL interpreter generates a sequence of discrete \texttt{Action} objects that must be executed within this real-time game loop.
Many user-facing operations---such as player movement, item manipulation, and inventory navigation---require multiple ticks to complete, necessitating a sophisticated action execution framework that bridges the gap between imperative DSL semantics and Minecraft's tick-driven engine.

\subsection{Tick-based Action Execution Model}

The action execution system is built around a queue-based architecture that processes actions one at a time, with each action given an opportunity to update its state on every game tick.

The core component is the \texttt{ActionRunner}, which is registered to receive callbacks from Minecraft's client tick events.
On each tick, the \texttt{ActionRunner} maintains a queue of pending \texttt{Action} objects and a reference to the currently active action (if any).
The tick method follows this pattern:

\begin{lstlisting}[language=Java]
public void tick() {
    if (current == null) {
        current = queue.poll();
        if (current == null) return;
    }

    switch (current.tick(ctx)) {
        case DONE    -> current = null;
        case RUNNING -> {}
    }
}
\end{lstlisting}

Each \texttt{Action} implements a \texttt{tick(MinecraftContext ctx)} method that returns an \texttt{ActionStatus}: either \texttt{RUNNING} (indicating that the action is incomplete and will continue on the next tick) or \texttt{DONE} (indicating that the action has completed and the next queued action should begin).

\subsection{Multi-tick Actions and Stateful Execution}

Many operations require coordination across multiple game ticks.
For example, moving the player forward by a specified distance requires sustained input over many frames, whereas waiting for a fixed duration requires tracking elapsed time across ticks.

\subsubsection{Waiting and Delay Actions}

The \texttt{WaitTicksAction} is one of the simplest multi-tick actions: it simply counts elapsed ticks until the specified duration is reached.
On each tick, it increments an elapsed counter and returns \texttt{RUNNING} until the counter reaches the target.
This enables DSL constructs such as \texttt{wait(60)} to insert precise delays between other actions:

\begin{lstlisting}[language=Java]
public class WaitTicksAction implements Action {
    private final int ticks;
    private int elapsed = 0;

    @Override
    public ActionStatus tick(MinecraftContext ctx) {
        return ++elapsed >= ticks ? ActionStatus.DONE : ActionStatus.RUNNING;
    }
}
\end{lstlisting}

Regarding player movement actions, the \texttt{MoveForwardAction} demonstrates the complexity of bridging user-level commands with real-time game physics.
When a DSL program calls \texttt{move\_forward(5)} to move forward five blocks, the action must:
\begin{itemize}
\item[--] Compute a target position based on the player's current facing direction.
\item[--] Check that the path from the player's current position to the target is walkable (accounting for terrain, obstacles, and step heights).
\item[--] Maintain simulated input (by forcing the forward key via \texttt{InputOverride}) across multiple ticks, allowing Minecraft's physics engine to move the player naturally.
\item[--] Monitor the player's actual position each tick to ensure progress is being made toward the target.
\item[--] Detect and handle obstacles or stuck conditions, terminating early if the player cannot progress.
\item[--] Release the forced input once the player has arrived at the target, returning \texttt{DONE}.
\end{itemize}

A simplified view of the state machine is:

\begin{lstlisting}[language=Java]
public ActionStatus tick(MinecraftContext ctx) {
    if (!initialized) {
        initialized = true;
        initializeTarget(ctx);
        if (!isPathClear(...)) return ActionStatus.DONE;  // Can't proceed
        InputOverride.INSTANCE.force(InputOverride.Key.FORWARD);
        forcedForward = true;
    }

    double distanceSq = distance(player, target);

    if (distanceSq <= ARRIVAL_DISTANCE_SQ) {
        InputOverride.INSTANCE.release(InputOverride.Key.FORWARD);
        return ActionStatus.DONE;
    }

    if (detectStuck(...)) {
        return ActionStatus.DONE;  // Fail-safe
    }

    return ActionStatus.RUNNING;
}
\end{lstlisting}

Crucially, the action does not directly move the player; instead, it manipulates Minecraft's input system (via \texttt{InputOverride}) and allows the game engine to apply physics and collisions naturally.
This ensures that movement respects the game's terrain, gravity, and collision detection without reimplementing Minecraft's movement logic.

\subsection{Input Override and Single-Tick Actions}

Many DSL operations can be completed in a single tick, such as selecting a hotbar slot, opening the inventory, or turning to face a cardinal direction.
However, these single-tick actions must not conflict with persistent input states from long-running actions like movement.

The \texttt{InputOverride} class manages a set of forced input keys (forward, backward, left, right, jump, sneak) using an \texttt{EnumSet}.
When an action forces a key \\ (e.g., \texttt{InputOverride.INSTANCE.force(InputOverride.Key.FORWARD)}), a mixin in the client input system checks this set before processing the player's actual keyboard input.
If a key is forced, the forced state takes precedence; otherwise, the player's actual input is used.

This design ensures that:
\begin{itemize}
\item[--] Long-running actions like movement can maintain continuous input across ticks without requiring the player to hold down keys.
\item[--] Single-tick actions (and the main game loop) can detect which keys are forced and respect them, avoiding conflicts.
\item[--] When an action completes, it can release its forced keys, restoring normal player control or allowing the next action to force different keys.
\end{itemize}

\subsection{Action Sequencing and Execution Flow}

The DSL interpreter produces an ordered sequence of \texttt{Action} objects, which are passed to the \texttt{ActionRunner} via an \texttt{ActionSequence}.
The \texttt{ActionSequence} is a builder-like container that holds an immutable list of actions:

\begin{lstlisting}[language=Java]
ActionSequence sequence = ActionSequence.start(action1)
    .then(action2)
    .then(action3);
runner.run(sequence);
\end{lstlisting}

When \texttt{runner.run(sequence)} is called, the \texttt{ActionRunner} replaces any existing queue and action with the new sequence's actions.
On each subsequent game tick, the runner calls the current action's \texttt{tick()} method.
When an action returns \texttt{DONE}, the runner polls the next action from the queue.
If the queue becomes empty, the runner remains idle until a new sequence is submitted.

This design ensures:
\begin{itemize}
\item[--] Actions are executed strictly sequentially, with each action completing before the next begins.
\item[--] Multi-tick actions automatically serialize with following actions; a movement command that takes 50 ticks will not overlap with subsequent inventory operations.
\item[--] The main game loop is never blocked; actions yield control every tick, allowing other game systems (rendering, entity updates, multiplayer synchronization) to proceed normally.
\end{itemize}

\subsection{Cleanup and Error Recovery}

The \texttt{stop} command and exception handling ensure that interrupted action sequences clean up their state properly.
When an exception occurs during interpretation or action execution, or when the player manually stops execution:

\begin{itemize}
\item[--] The \texttt{ActionRunner} is canceled, clearing the current action and queue.
\item[--] All forced input keys are released via \texttt{InputOverride.INSTANCE.releaseAll()}, restoring normal player control.
\item[--] The default keyboard input handler is restored via \texttt{MinecraftContext.restoreDefaultInput()}, undoing any input overrides and ensuring player-controlled movement is possible.
\item[--] Any open inventory screens are left in their current state (to avoid abrupt UI changes that might confuse players), but subsequent player input is once again processed normally.
\end{itemize}

This recovery mechanism ensures that even in edge cases, the player regains full control of the game and the Minecraft client remains stable.

\subsection{Design Rationale}

The tick-based action model was chosen because it aligns naturally with Minecraft's game loop and provides several benefits:

\begin{itemize}
\item[--] It provides Real-Time Interaction as the player can continue to observe and interact with the world while DSL actions execute, making the system feel responsive rather than blocking.

\item[--] The model is robust by relying on Minecraft's built-in input and physics systems rather than attempting to reimplement them, the system is robust to game updates, mod interactions, and edge cases in terrain or entity behavior.

\item[--] There is an established simplicity since each action class is stateful and self-contained, making it easy to reason about and test in isolation. Complex multi-tick behavior emerges from simple single-tick state machines.

\item[--] The system is extensible, since new actions can be added by implementing the \texttt{Action} interface and registering them in the standard library, without modifying the core runner or tick loop.

\item[--] There is also an aspect of observability, as players can enable debug logging to watch each action transition from \texttt{RUNNING} to \texttt{DONE}, aiding in learning and troubleshooting.
\end{itemize}

